\documentclass{article}
\usepackage{amsmath}
\usepackage{graphicx}
\usepackage{fontspec}  % compile with XeLaTeX (for the rupee sign)
\title{Solar Power for Village Schools}
\author{Format X Test Team}
\begin{document}
\maketitle

\section*{1. Why Schools Need Solar Power}

Many village schools lose electricity for several hours every day.
Without power, fans stop, lights go out, and computers cannot be used.
Teachers often move classes outside when the rooms become too dark.
Students lose valuable learning time during these power cuts.
Diesel generators can help, but fuel is expensive and noisy.
Solar panels offer a cleaner and cheaper way to keep schools running.
A small rooftop system can power lights, fans, and a few computers.
This report studies one such system over five months.

The school in this study has 240 students and 9 teachers.
Before the project, the school had power for only 14 hours a day on average.
The local R\&D team designed the system with help from two engineers.
The project budget was ₹1,20,000, including panels, batteries, and wiring.
A similar imported kit was quoted at about \$450 per panel.
All readings were saved in a file called solar\_data\_2026.csv for later study.
The results are summarised in the sections below.

\section*{2. How Much Energy a Panel Makes}

The energy a solar panel produces depends on four main things.
The first is the area of the panel, measured in square metres.
The second is the panel's efficiency, which is the share of sunlight it turns into electricity.
The third is the number of peak sun hours in a day.
The fourth is the performance ratio, which accounts for losses in wires, heat, and dust.
Engineers combine these four values in a simple formula.
The formula below gives the energy produced in one day.
Each symbol is explained right after the formula.

\begin{equation}
E = A \times r \times H \times PR
\label{eq:1}
\end{equation}

In this formula, E is the daily energy in kilowatt-hours.
A is the total panel area, and r is the panel efficiency.
H is the number of peak sun hours, and PR is the performance ratio.
The school's panels cover 20 square metres with an efficiency of 18\%.
The region gets about 5 peak sun hours on a typical day.
With a performance ratio of 0.75, the system should make about 13.5 kWh per day.
This is enough to run the school's lights, fans, and computer lab.

\section*{3. Measured Results}

The system was switched on in January and measured every day.
A small meter recorded the energy produced each hour.
Sunny days were counted using a simple light sensor.
The table below shows the results for the first five months.
Energy is given in kilowatt-hours, and savings are given in rupees.
Savings were calculated at a rate of ₹8 per unit of electricity.
January and March gave the best results because of clear skies.
May was lower because of early rain and more cloudy days.

\begin{table}[h]
\centering
\label{tab:results}
\begin{tabular}{lrrr}
Month & Sunny days & Energy (kWh) & Savings (₹) \\
\hline
January & 24 & 405 & 3,240 \\
February & 23 & 372 & 2,976 \\
March & 26 & 418 & 3,344 \\
April & 22 & 381 & 3,048 \\
May & 18 & 330 & 2,640 \\
\end{tabular}
\end{table}

In total, the system produced 1,906 kWh in five months.
This saved the school ₹15,248 on its electricity bill.
The best single day produced 15.2 kWh, which is more than the formula predicted.
The worst day produced only 3.1 kWh during a heavy storm.
On average, the system made 12.6 kWh per day.
This is about 93\% of the value predicted by the formula.
The small gap is mostly explained by dust on the panels.

\section*{4. The Installation}

The panels were mounted on the flat roof of the main building.
They face south and are tilted at 25 degrees to catch the most sunlight.
A battery bank in the store room keeps power available after sunset.
The wiring runs through a safety switch and a small control box.
Local workers finished the installation in four days.
Two teachers were trained to read the meter and clean the panels.
The photo below shows the finished system on the roof.

\begin{figure}[h]
\centering
\includegraphics[width=0.8\textwidth]{images/solar_panels.png}
\caption{Solar panels on the roof of the main school building}
\label{fig:roof}
\end{figure}

The panels are placed in four rows with space for cleaning between them.
The roof was checked by an engineer before the work began.
No leaks or cracks were found during the first five months.
Birds sometimes sit on the frames, so a simple net was added.
The panels are cleaned with water every two weeks.
Cleaning takes less than one hour for two people.

\section*{5. Costs and Savings}

The most important question for the school was when the system would pay for itself.
This is called the payback period.
It is found by dividing the total cost by the yearly savings.
The yearly energy was estimated from the five-month results.
The formula below shows the calculation.
Here, C is the total cost, E\_year is the yearly energy, and p is the price of one unit.

\begin{equation}
T = \frac{C}{E_{year} \times p}
\label{eq:2}
\end{equation}

The system is expected to make about 4,600 kWh in a full year.
At ₹8 per unit, this saves about ₹36,800 every year.
With a total cost of ₹1,20,000, the payback period is about 3.3 years.
After that, the electricity is almost free for the rest of the panels' 25-year life.

\section*{6. Conclusion}

The solar system gave the school reliable power during school hours.
Students no longer lose class time because of power cuts.
The measured energy matched the formula within 7\%.
The money saved can now be spent on books and science kits.
Other schools can use the same formula to plan their own systems.
The main lessons are to choose a strong roof and to keep the panels clean.
With simple care, a small solar system can serve a school for many years.

\end{document}
